\section*{الفصل \ref{ch:b3:pericyclic} --- التفاعلات المحيطية الحلقية}

\begin{solution}{exo:b3:pericyclic:1}
ديلز--ألدر: \omterm{def:b3:pericyclic:families}{إضافة حلقية} [4+2]. انفتاح السيكلوبيوتين: حلقي كهربائي. كوب: \omterm{def:b3:pericyclic:families}{إعادة ترتيب سيغماتروبية}
[3,3]. فقد السلفولين للجزيء $\ce{SO2}$: خيليتروبي. الأوزون والألكين: \omterm{def:b3:pericyclic:dipolar}{إضافة حلقية ثنائية القطب 1،3}
([3+2]). الانزياح \babelsublr{$[1,5]$-H}: سيغماتروبي.
\end{solution}

\begin{solution}{exo:b3:pericyclic:2}
الهكساترايين (6 إلكترونات): \omterm{def:b3:pericyclic:rotation-modes}{متعاكس الدوران} حراريًا، \omterm{def:b3:pericyclic:rotation-modes}{ومتساير الدوران} ضوئيًا.
البوتاديين (4): \omterm{def:b3:pericyclic:rotation-modes}{متساير الدوران} حراريًا، \omterm{def:b3:pericyclic:rotation-modes}{ومتعاكس الدوران} ضوئيًا.
\end{solution}

\begin{solution}{exo:b3:pericyclic:3}
أربعة إلكترونات، حراريًا: انفتاح \omterm{def:b3:pericyclic:rotation-modes}{متساير الدوران}. ومجموعتا الميثيل، على وجهين متقابلين
للحلقة، تدوران كلتاهما إلى الخارج: (\textit{E},\textit{E})-هكسا-2،4-ديين.
\end{solution}

\begin{solution}{exo:b3:pericyclic:4}
التفاعل [2+2] \omterm{def:b3:pericyclic:faciality}{الأحادي الوجه} فيه $4q$ إلكترونًا: ففي الحالة الأساسية يتداخل HOMO أحد
الإيثينين مع LUMO الآخر بطرف رابط وطرف مضاد للترابط.
وفي الحالة المثارة يؤدي المدار $\pi^*$ المشغول بإلكترون واحد في أحد الجزيئين دور
HOMO، فيتوافق الطرفان كلاهما.
\end{solution}

\begin{solution}{exo:b3:pericyclic:5}
يُحفظ \omterm{def:b3:point-groups:mirror}{مستوي المرآة}. الهكساترايين: $\psi_1$ S، $\psi_2$ A، $\psi_3$ S مشغولة؛ $\psi_4$ A، $\psi_5$ S، $\psi_6$ A فارغة.
السيكلوهكساديين، بالترتيب: $\sigma$ S، $\pi_1$ S، $\pi_2$ A مشغولة؛ $\pi_3^*$ S، $\pi_4^*$ A، $\sigma^*$ A فارغة. وبالوصل بحسب
التناظر بالترتيب، $\psi_1 \to \sigma$، $\psi_2 \to \pi_2$، $\psi_3 \to \pi_1$: مشغول إلى مشغول، فالانغلاق \omterm{def:b3:pericyclic:rotation-modes}{المتعاكس الدوران}
مسموح حراريًا.
\end{solution}

\begin{solution}{exo:b3:pericyclic:6}
ستة إلكترونات تحت الضوء: \omterm{def:b3:pericyclic:rotation-modes}{متساير الدوران}، فتنتهي مجموعتا الميثيل الطرفيتان على
وجهين متقابلين: \textit{trans}-5،6-ثنائي ميثيل سيكلوهكسا-1،3-ديين (والتفاعل الحراري يعطي
المتماكب \textit{cis}).
\end{solution}

\begin{solution}{exo:b3:pericyclic:7}
ينتقل الهيدروجين الذي على C5 إلى C1 في الديين، جاره حول الحلقة،
عبر حالة انتقالية سداسية الإلكترونات ($\sigma2_s + \pi4_s$)، \omterm{def:b3:pericyclic:faciality}{على نحو أحادي الوجه}: فهجرة واحدة
تعطي 1-ميثيل سيكلوبنتاديين، وثانية 2-ميثيل سيكلوبنتاديين. والسلسلة الخماسية الكربون تسمح
للهيدروجين بالبقاء على وجه واحد، وهو ما تفرضه الحلقة الصغيرة على أي حال.
\end{solution}

\begin{solution}{exo:b3:pericyclic:8}
$\pi2_s$ (الألكين) $+ \pi2_a$ (الرابطة C=C في الكيتين، المهاجَمة على وجهين متقابلين عبر
المدار $\pi^*$ العمودي للرابطة C=O): عدد المكوّنات $(4q + 2)_s$ واحد (الألكين)، فردي:
مسموح حراريًا.
\end{solution}

\begin{solution}{exo:b3:pericyclic:9}
في الكرسي، الرابطة C=C في مجموعة الفينيل (الذرتان 1--2)، والأكسجين (3)، وكربون
$\ce{OCH2}$ (4) والرابطة C=C في الأليل (5--6)؛ تنكسر الرابطة O--C وتتكوّن الرابطة C1--C6:
والناتج بنت-4-إينال،
$\ce{OHC-CH2-CH2-CH=CH2}$.
\end{solution}

\begin{solution}{exo:b3:pericyclic:10}
ثمانية إلكترونات هي $4q$؛ ومع مكوّن واحد \omterm{def:b3:pericyclic:faciality}{متعاكس الوجهين} يكون الترتيب من طوبولوجيا
موبيوس، العطرية مع $4q$ إلكترونًا: \omterm{def:b3:pericyclic:mobius}{فالحالة الانتقالية عطرية} 
والتفاعل مسموح حراريًا (كالانغلاق \omterm{def:b3:pericyclic:rotation-modes}{المتساير الدوران} للأوكتاتترايين).
\end{solution}

\begin{solution}{exo:b3:pericyclic:11}
مزاوجة أكبر المعاملات تصل N3 بالكربون الطرفي: فيتكوّن
الترايازول الثنائي الاستبدال 1،4. ودون حفاز تكاد المعاملات لا تختلف 
ويتكوّن المتماكبان 1،4 و1،5 كلاهما؛ والنحاس(I) يجعل التفاعل متدرجًا عبر
أسيتيليد النحاس ولا يعطي إلا المتماكب 1،4.
\end{solution}

\begin{solution}{exo:b3:pericyclic:12}
يجسر الهيدروجين بين طرفي SOMO لجذر من $j$ ذرة، $k = (j + 1)/2$. ومعاملاه
الطرفيان متناسبان مع $\sin(k\pi/(j + 1))$ و$(-1)^{k+1}\sin(k\pi/(j + 1))$. ومن أجل $j = 7$، $k = 4$: إشارتان متعاكستان، ومنه
تكون الفصوص ذات الإشارة نفسها على وجهين متقابلين عند الطرفين: \omterm{def:b3:pericyclic:faciality}{متعاكس الوجهين}. ومن أجل
$j = 5$، $k = 3$: الإشارة نفسها على الوجه نفسه: \omterm{def:b3:pericyclic:faciality}{أحادي الوجه}.
\end{solution}

\begin{solution}{pb:b3:pericyclic:1}
\textbf{1.} الرابطة $\sigma$ بين C9 وC10 في الحلقة B؛ ومع الديين 5،7 تعطي هكساترايينًا،
C10=C5--C6=C7--C8=C9.

\textbf{2.} ستة إلكترونات (رابطتا $\pi$ والرابطة $\sigma$): انفتاح حلقي كهربائي.

\textbf{3.} حراريًا، لا يتوازن السيكلوهكساديين والهكساترايين المفتوح المقابل له إلا
ببطء، عبر حاجز مرتفع، والحلقة هي الطرف الأكثر استقرارًا؛ أما
الديين فيمتص الضوء فوق البنفسجي، وتنفتح حالته المثارة في غضون
بيكوثانيات.

\textbf{4.} \omterm{def:b3:pericyclic:rotation-modes}{متساير الدوران}: ففي الحالة المثارة يكون للمدار LUMO المشغول بإلكترون واحد في
نظام الترايين ($k = 4$) طرفان متعاكسا الإشارة.

\textbf{5.} تأتي الرابطة C6=C7 من الحلقة: فامتدادا السلسلة عندها، C5 وC8،
كانا على الجانب نفسه منها، ويبقيان كذلك: \textit{Z}.

\textbf{6.} الانفتاح الحراري، المسموح \omterm{def:b3:pericyclic:rotation-modes}{بالدوران المتعاكس} وحده، له حاجز أبعد بكثير مما
تستطيع حرارة الجسم توفيره، وسيقود صعودًا، نحو الترايين الأقل استقرارًا.

\textbf{7.} من C19 (الميثيل على C10) إلى C9: انزياح سيغماتروبي \babelsublr{$[1,7]$-H}.

\textbf{8.} ثمانية: إلكترونات $\pi$ الستة في الترايين وإلكترونا الرابطة C--H.

\textbf{9.} $\sigma2 + \pi6$. فإذا كانا أحاديي الوجه كلاهما لكان كل منهما s، مع مكوّن $(4q + 2)_s$ واحد من كل منهما:
زوجي، ممنوع. ومع ترايين \omterm{def:b3:pericyclic:faciality}{متعاكس الوجهين}، $\sigma2_s + \pi6_a$: مكوّن واحد معدود، فردي،
مسموح.

\textbf{10.} سبع ذرات في سلسلة حلزونية ذات تشكيل \textit{Z} تتيح للهيدروجين أن ينتقل من وجه إلى
الآخر؛ أما السلسلة الثلاثية الذرات فلا تستطيع أن تلتوي إلى هذا الحد.

\textbf{11.} موبيوس (مكوّن واحد \omterm{def:b3:pericyclic:faciality}{متعاكس الوجهين}) مع ثمانية إلكترونات، $4q$: عطرية.

\textbf{12.} إنه مسموح حراريًا؛ ويُعبر حاجزه ببطء عند درجة حرارة الجسم.

\textbf{13.} ستة إلكترونات تحت الضوء: \omterm{def:b3:pericyclic:rotation-modes}{متساير الدوران} من جديد، لكنه يستطيع أن يدير الطرفين في
الاتجاه الآخر فيضع H9 وميثيل C19 على الوجهين الآخرين: فيعطي متماكبًا فراغيًا
مغلق الحلقة، هو اللوميستيرول.

\textbf{14.} مع رابطة مركزية \textit{E}، يكون C19 وC9 على جانبين متقابلين من السلسلة وبعيدين
أحدهما عن الآخر: فلا يستطيع الهيدروجين البلوغ.

\textbf{15.} تمتص طليعة الفيتامين \babelsublr{D$_3$} نفسها الضوء فتتحول إلى لوميستيرول وتاكيستيرول،
اللذين لا يعطيان الفيتامين D؛ فيُبلغ مزيج مستقر ضوئيًا، يضع سقفًا
للإنتاج.

\textbf{16.} يحدث الانفتاح في الحالة المثارة، في فيمتوثوانٍ إلى بيكوثوانٍ؛ أما
الانزياح فتفاعل حراري بحاجز قدره \qty{85}{kJ/mol}.

\textbf{17.} $t_{1/2} = \ln 2/k = 0.693/\qty{0.023}{h^{-1}} = \qty{30}{h}$.

\textbf{18.} $1 - \eu^{-0.023 \times 24} = 0.42$.

\textbf{19.} $\ln 10/k = 2.303/\qty{0.023}{h^{-1}} = \qty{100}{h}$.

\textbf{20.} $k(\qty{293.15}{K}) = k(\qty{310.15}{K})\,\eu^{-(E_a/R)(1/293.15 - 1/310.15)} = 0.023 \times \eu^{-1.91} = \qty{3.4e-3}{h^{-1}}$.

\textbf{21.} $2.303/\qty{3.4e-3}{h^{-1}} = \qty{680}{h}$، أي نحو 28 يومًا.

\textbf{22.} الخطوة الحرارية أسرع بسبع مرات عند \qty{37}{\celsius} منها عند \qty{20}{\celsius}، بحيث تتحول
طليعة الفيتامين المصنوعة في صباح مشمس في غضون أيام.

\textbf{23.} الهيكل الستيرويدي، بمراكزه الفراغية الكثيرة، يأتي جاهزًا من
الستيرول؛ ثم يُجري الضوء والدفء الخطوتين المحيطيتين الحلقيتين تمامًا كما في
الجلد.

\textbf{24.} عند درجة حرارة الجسم، تصير 90\,\% من طليعة الفيتامين \babelsublr{D$_3$} فيتامين \babelsublr{D$_3$} في نحو
\qty{100}{h}، أي أربعة أيام.
\end{solution}
